\section{Introduction} %
\label{sec:intro}

% BD: trying to get away from talking about verification 'procedures'
% here -- it may prime the reader for a type checking algorithm

Static type systems are one of the most versatile and widely-used
program reasoning techniques available today. While their particulars
can vary wildly, common to all type systems are two key ingredients:
the use of \emph{types} to specify program behavior, and some
deductive mechanism for \emph{compositionally verifying} that a
program is consistent with a specification. Type systems have
traditionally targeted \emph{safety} properties, a focus embodied by
Milner's famous slogan that ``well-typed expressions cannot 'go
wrong'''~\cite{Milner+78}. Here, the typing judgement $\Gamma \vdash e
: \tau$ asserts that $e$ \emph{may} evaluate to many values (under an
environment consistent with $\Gamma$), but that set of values is
entirely contained within the ``safe'' region specified by $\tau$.

Dually, \emph{reachability} properties seek to precisely capture a
subset of the values that a program \emph{must} reach. Rather than
ruling out certain executions, reachability analyses seek to soundly
establish the existence of feasible executions.  Reachability analyses
are a well-established and practical technique for informing program
optimizations and transformations~\cite{WS94,GM21,JTW+98}.  They have
also been effectively used to find buggy behaviors~\cite{BMC+03}; for
example, several recent works~\cite{GOS19,ReverseHoareLogic, OHP19,
  ZDS23} have studied ``incorrectness'' systems designed to prove the
existence of witnesses of buggy program behavior, with a goal of
avoiding the overwhelming number of false positives reported by
typical bug detection tools.  The question of how to develop type
systems targeting reachability properties, however, has been largely
unexplored. A few recent exceptions~\cite{RefinementTypeRefutations,
  CoverageType, Ill-types} have studied type-based reachability
reasoning, but each of these systems has focused on a particular
property, e.g., checking coverage properties of test input
generators~\cite{CoverageType}, or applications, e.g., generating
concrete counterexamples refuting refinement type
specifications~\cite{RefinementTypeRefutations} or proving ill-typed
programs do not evaluate~\cite{Ill-types}, rather than providing a
general type-based reachability framework.

% BD: One thing we should probably foreground more is that our
% semantic model can track dependencies between free variables when
% reasoning compositionally...
This paper proposes a novel semantic model that can simultaneously
capture both over- and under-approximations of program behaviors,
allowing safety and reachability specifications to be compositionally
combined within a single type discipline. This model gives rise to a
type system in which type qualifiers are equipped with a modality
indicating whether it should be interpreted as an over- or
under-approximate specification. This type system naturally supports
higher-order, deterministic, and nondeterministic programs, captures
both may- and must-style function summaries~\cite{GN+10}, and can be
applied to a range of verification tasks, including
incorrectness-style reasoning about buggy programs~\cite{OHP19} and
proving completeness of test-input generators~\cite{CoverageType}.

To illustrate, consider the implementation of Peano division from
Rocq's standard library:
\begin{align}
  \lambda x.\lambda y. \zif{y \leq 0}{x}{\Code{fst~(divmod}~x~(y - 1)~0~(y - 1))}
  \tag{\Code{/}}\label{ex:div}
\end{align}
A client of \ref{ex:div} may want to know both that (a) on positive
inputs, its result is always a positive number no greater than its first
argument (a safety property); and (b) for every positive number $n$,
there exist positive inputs $x$ and $y$ such that $x$ \ref{ex:div}
$y = n$ (a reachability property). Conventional refinement type
systems, which interpret refinements as overapproximations, can verify
the former, but not the latter: e.g., a function that always returns $1$
also satisfies (a). Conversely, an underapproximate interpretation can
establish (b), but would also admit an implementation that always
returns $x$. Our type system, in contrast,  captures both kinds of
properties as types:
\begin{align}
  & x{:}\ort{\Nat}{\vnu > 0} \sarr y{:}\ort{\Nat}{\vnu > 0 } \sarr \ort{\Nat}{x \geq \vnu > 0} \tag{$\tau_a$}\label{ex:tauA}\\
  & x{:}\urt{\Nat}{\vnu > 0} \sarr y{:}\urt{\Nat}{\vnu > 0} \sarr \urt{\Nat}{\vnu > 0 } \tag{$\tau_b$}\label{ex:tauB}
\end{align}\noindent
% \BD{One of the tricky things here is that \ref{ex:div} does not have
%   the bottom signature -- we'd need to use bunched implication to
%   capture the property, and we do not want to introduce that just
%   yet.}
The result type of \ref{ex:tauA} uses $\ort{b}{\phi}$ to express that
the qualifier $\phi$ should be interpreted as a safety constraint: a
term with this type cannot produce any value outside the set of values
the type denotes.  Conversely, \ref{ex:tauB} uses $\urt{b}{\phi}$ to
capture a reachability constraint: a term with this type must produce
every value found in the set the type denotes.

The intuitive meaning of each modality changes subtly when a
refinement appears in a typing context instead of a result type. To
establish the typing judgment $x{:} \ort{b}{\phi} \vdash e ~:~\tau$,
our type system intuitively treats any free occurrences of the
variable $x$ \emph{demonically}, since the surrounding environment is
free to choose any value of $x$ consistent with $\phi$, which may
cause $e$ to violate the constraints defined by $\tau$. Conversely,
when establishing the typing judgment $x{:} \urt{b}{\phi} \vdash e
~:~\tau$, $x$ is treated \emph{angelically} since it must be able to
take on every value consistent with $\phi$.

The availability of first-class functions adds additional
complications. Higher-order functions like the following, for example,
naturally expose two distinct contexts:
$x : \tau_x \vdash \lambda y. \ldots ~:~ \tau_y \sarr \tau$.  The
first, $x{:} \tau_x$, captures the type of the free variable $x$,
while the second, $y{:} \tau_y$, provides assumptions about the
function's arguments at the point it is applied. Managing these
contexts separately is necessary because the instantiation of the free
variables in a function and its parameters can occur at different
points in a program's evaluation. To illustrate, consider the
following (failing) type judgement:
\begin{align*}
    x{:}\urt{\Nat}{\vnu > 0} &\vdash  \lambda f. f\;x ~~:~~ (\urt{\Nat}{\vnu > 0}\sarr\tau_f) \sarr \urt{\Nat}{\vnu > 0} \quad\judgfail
\end{align*}\noindent
The higher-order function being typed takes a function $f$ and
directly applies it to $x$. The type binding for $x$ in the type
context assumes it can reach every positive number.  Without
constraining the type of $f$ to account for the free variable $x$ in
the expression's closure, we cannot guarantee that an arbitrary
function supplied for $f$ will satisfy the reachability requirement of
this expression's result.  For example, applying this expression to
the function $\lambda y.x - y$ yields the constant function that
always returns $0$.
\begin{align*}
  & \zlet{x}{e_x}{\boxed{(\lambda f. f\;x)\;(\lambda y.x - y)}} \;\hookrightarrow^*
    \; \zlet{x}{e_x}{\boxed{x - x}} \;\hookrightarrow^*\; 0
\end{align*}\noindent
The problem is that the ``captured context'' of function
$\lambda y.x - y$ contains variable $x$ \emph{and} that this function
is supplied $x$ as its argument.  Thus, for this type judgement to be
valid, the type of $f$ must enforce that its parameter $y$ can be
chosen independently of $x$. This requires our type system to reason
explicitly about dependencies between a function's parameters and
captured variables when both are treated angelically.

The previous example showed how dependencies between captured
variables and function parameters can break compositional reachability
reasoning.  The example given below illustrates a more subtle problem
that can arise when the result of an intermediate function application
depends on a captured variable, constraining the angelic choices
available to subsequent applications in a way that also breaks
compositional reasoning:
\begin{align*}
    x{:}\urt{\Nat}{\vnu > 0}, y{:}\urt{\Nat}{\vnu > 0} \vdash \lambda f. f(f\;y) : \tau_f \sarr \urt{\Nat}{\vnu > 0} \quad \judgfail
\end{align*}\noindent
A counter-example to the validity of this judgement is given below:
\begin{align*}
    &g \doteq \lambda y.\zif{y = 1}{x}{2} \quad
    \\& \zlet{x}{e_x}{\zlet{y}{e_y}{\boxed{(\lambda f. f\;(f\;y))\;g}}} \;\hookrightarrow^*\;
    \\& \zlet{x}{e_x}{\boxed{g\;(g\;1)}} \;\hookrightarrow^*\;  \quad
  \text{\textcolor{gray}{(we assume angelically that $e_y \hookrightarrow^* 1$)}}
    \\& \zlet{x}{e_x}{\boxed{g\;x}} \;\hookrightarrow^*\; \zlet{x}{e_x}{\boxed{\zif{x = 1}{x}{2}}} \;\hookrightarrow^*\; \boxed{1 \text{ or } 2}
\end{align*}\noindent
Angelically choosing $y = 1$, the first application ($g\; 1$) returns
$x$. This result is then passed to $g$ again, but since $x$ is a
captured variable rather than a fresh angelic choice, the second
application $(g\; x)$ can only return 1 or 2, and so the nested
application fails to reach any positive number greater than 2. Thus,
although the individual applications of $g$ appear to satisfy its
reachability specification, their composition does not, again
invalidating compositional reasoning.

\iffalse
% Old Intro
A typical type-based verification procedure provides a compositional
way to reason about a program's \emph{safety} properties~\cite{Milner+78}.
The procedure can be thought of as defining a logical deduction system
that validates a program's behaviors with respect to its type
specification.  For example, a typing judgement $\Gamma \vdash e :
\tau$ asserts that $e$ \emph{satisfies} the specification given by type
$\tau$ under a typing context $\Gamma$.  This \emph{demonic}
interpretation provides a safety guarantee: the result of a well-typed
term is always contained in $\tau$.  Concretely, the evaluation of $e$
\emph{may} produce a value in $\tau$'s denotation (but nothing else),
and must evaluate in \emph{all} environments whose bindings are
consistent with $\Gamma$.

A less studied, but equally valuable, verification procedure is one in
which types are used to establish \emph{reachability}, rather than
safety, properties.  Here, a typing judgement asserts that $e$
\emph{realizes} the specification given by type $\tau$ under a typing
context $\Gamma$.  This \emph{angelic} interpretation provides a
reachability guarantee: every value that inhabits $\tau$ must be the
result of a well-typed term.  Concretely, the evaluation of $e$
\emph{must} produce every value in $\tau$'s denotation (but possibly
more), and must evaluate in \emph{some} environment whose bindings are
consistent with $\Gamma$.

Reachability reasoning~\cite{ReverseHoareLogic, OHP19, ZDS23} has
typically been used to prove claims about ``incorrectness'' properties
of programs, i.e., the existence of witnesses of buggy program
behavior that can avoid overwhelming false positives in bug detection.
There have been relatively few instances that apply this style of
reasoning in type-based verification, however. Recent refinement type
systems~\cite{RefinementTypeSymbolicExecution,RefinementTypeRefutations}
instantiate must-style reasoning to refute type specifications by
searching for concrete counterexamples.  Other type systems focus on
specific reachability properties: for example, \emph{coverage
  types}~\cite{CoverageType} ensure that all desirable test inputs
will be produced by a test input generator, while~\citet{Ill-types}
guarantees that all ill-typed programs will definitely get stuck.

% Existing approaches, however, typically fix on one style of reasoning
% or the other and are not equipped to simultaneously consider both
% safety and reachability properties together, limiting their ability to
% address useful verification tasks.  For example, consider a
% specification for subtraction ($\Code{-}_{\Code{natpos}}$) that types
% argument $y$ as a natural number less than
% $x$ and asserts that the result is a natural number greater
% than 0 but at most $x$:
% \[x : \{\Code{nat}\, |\ 0 < \nu\}, y : \Code{\{\Code{nat}\, |\, 0 \leq \nu < x\}}
% \vdash x\;\Code{-}_{\Code{natpos}}\;y : \{\Code{nat}\, |\, 0 < \nu
% \leq x \}\] This type has both a safety component --
% $\Code{-}_{\Code{natpos}}$ must produce a result in the range (0,$x$];
%   and a reachability component -- every natural number in (0,$x$] can
%     be produced by instantiating $y$ appropriately.  By considering
%     just safety, we could verify the former but would be unable to
%     guarantee reachability.  For example, an implementation that
%     simply returned the constant 1 would be vacuously safe, but would
%     not satisfy reachability.    Conversely, by
%     viewing this specification from a reachability lens, we could
%     verify the latter but would be unable to guarantee that the
%     operation is actually safe.  For example, an implementation that
%     returns an arbitrary positive number satisfies the reachability
%     requirement since it can return every value in $(0, x]$, but is
%       unsafe since it may also return values outside this range.  A
%       unified treatment would simultaneously establish that the
%       implementation is both \emph{sound} and \emph{complete} with
%       respect to its specification.

Moreover, existing approaches typically fix on one style of reasoning
or the other and are not equipped to simultaneously consider both
safety and reachability properties together, limiting their ability to
address useful verification tasks.

As a simple example, consider a specification for subtraction over
natural numbers ($\Code{-}_{\Code{nat}}$) that types argument $x$ and
$y$ as non-negative integers and asserts that the result is a
non-negative integer less than or equal to $x$\footnote{In the
  following, we drop the ``\Code{nat}'' subscript when obvious from
  context.}:
  \[x : \{\Code{int}\, |\ 0 \leq \nu\}, y : \Code{\{\Code{int}\, |\, 0 \leq \nu \}}
      \vdash x\;\Code{-}_{\Code{nat}}\;y : \{\Code{int}\, |\, 0 \leq
      \nu \leq x \}\]
\noindent A standard implementation is $\Code{-}_{\Code{nat}} \doteq
\lambda x.\lambda y. \zif{x \leq y}{0}{x - y}$.  Its type has both (a)
a safety component -- the function $\Code{-}_{\Code{nat}}$ must
produce a result in the range $[0,x]$; and (b) a reachability
component -- every integer in $[0,x]$ can be produced by instantiating
$y$ appropriately.  By considering just safety, however, we could
verify the former but would be unable to guarantee reachability.  For
example, an alternative (incorrect) implementation that ignores its
inputs and simply returns the constant $1$ would be vacuously safe,
but would not capture the desired reachability property that every natural number
in the range $[0,$x$]$ can be expressed as the difference of two other numbers.  Conversely,
by viewing this specification from a reachability lens, we could
verify (b) but would be unable to guarantee that the operation is
actually safe.  For example, an alternative (incorrect) implementation
that ignores its inputs and simply returns the output of a random
number generator satisfies our reachability requirement since it can
potentially return every value in $[0, x]$, but is clearly unsafe
since it may also return values outside this range.  A unified
treatment would simultaneously establish that the standard
implementation is both \emph{sound} and \emph{complete} with respect
to its type specification.

Unifying safety and reachability reasoning in a type-based setting
is challenging in functional languages, where functions may have
free variables whose types determine whether a safety or reachability
guarantee can be established.  Consider the following example:
\begin{align*}
    &x{:}\urt{\Nat}{\vnu > 0} \vdash \lambda y.x - y : y{:}\urt{\Nat}{\vnu > 0} \sarr \urt{\Nat}{\vnu > 0} \quad\ \textcolor{green}{\text{\ding{51}}}
    \\&x{:}\ort{\Nat}{\vnu > 0} \vdash \lambda y.x - y : y{:}\urt{\Nat}{\vnu > 0} \sarr \urt{\Nat}{\vnu > 0} \quad\judgfail
\end{align*}\noindent
We use $\ort{b}{\phi}$ to express that a primitive type $b$ refined
with a qualifier $\phi$ should be interpreted as a safety constraint:
a term that has this type cannot produce any value outside the set of
values the type denotes.  Conversely, $\urt{b}{\phi}$ establishes a
reachability constraint: a term that has this type must produce every
value found in the set the type denotes.  Thus, the type
$y{:}\urt{\Nat}{\vnu > 0} \sarr \urt{\Nat}{\vnu > 0}$ specifies that
the function, $\lambda\, y. x - y$, must be able to produce every
positive number from its input $y$, whose type also guarantees that it
must produce every positive number.  The two judgements differ only in
the type ascribed to the free variable $x$.  In the first judgement,
$x$'s type mandates it produce every positive number.  This means that
for any desired output $v$, we can angelically pick values for $x$ and
$y$ consistent with their types such that their difference yields $v$.
For example, to produce output $v$, we can choose $x = v + 1$ and $y =
1$, both of which are positive and satisfy their respective types, and
whose difference is $v$.  In the second judgement, $x$'s type merely
constrains its value to be some positive number.  Since $x$'s type
requires that we demonically account for all positive values, no
single choice works for every combination of $y$ and output $v$, and
so the judgement is not valid.

\fi

% In order to define a type theory that uniformly supports both safety
% and reachability reasoning, we treat these two styles of
% interpretation as first-class modalities, and support fine-grained
% mixing in arbitrary ways.
To solve the context dependency issues raised above, we equip typing
contexts with \emph{capabilities} that are treated as \emph{resources}
that mediate demonic or angelic choices as necessary. We then
introduce a \emph{resource algebra} to compose and destruct these
capabilities in a principled way.  This approach enables us to handle
context disjointness analogously to the way bunched
implications~\cite{OHP99} handle resource separation, lifting the
additive/multiplicative conjunction distinction from the level of
propositions to the level of type contexts.

This framing naturally leads to the notion of a \emph{context type}
that unifies safety and reachability reasoning in a pure functional
language.  Specifically, we allow type contexts to be related via
dependent ($\Gamma_1 \ctcomma \Gamma_2$) and disjoint
($\Gamma_1 \ctstar \Gamma_2$) conjunction.  The type constructor
$\ctsarr$ internalizes dependent conjunction, typing functions whose
argument context is entangled with the closure context, while
$\ctwand$ internalizes independent conjunction, typing functions whose
argument context does not depend on the closure, directly reflecting
whether angelic choices for the argument are constrained by the
closure.

This bunched type system thus enables the tracking of complex
data-flow dependencies and guides the angelic choices that are
permissible in a typing judgement.  Specifically, since contexts can
be combined via context combinators, the structure of a context is
naturally a binary tree whose internal nodes record how sub-contexts
are combined, enabling frame-based local reasoning.  For example, to
establish that \ref{ex:div} can produce every positive number, we can
type it against the following type signature:
\begin{align*}
  & \vdash \ref{ex:div} ~:~ x{:}\urt{\Nat}{\vnu > 0} \sarr y{:}\urt{\Nat}{\vnu > 0} \ctwand \urt{\Nat}{\vnu > 0} \quad \judgok
\end{align*}\noindent
Here, the type constructor ($\ctwand$) requires the arguments for $x$
and $y$ to be independent, guaranteeing that each angelic choice can
be made regardless of the other.

%SJ: This sentence is a bit too cryptic - it teases non-determinism
% without really providing enough motivation.  I think it would be
% fine to list it in the contributions as we do, and have the reader
% wait until that section to get details.
%% While the result context of a
%% deterministic function is always dependent on its input context,
%% \emph{nondeterministic} functions may produce results that are
%% independent of some inputs, a distinction that our system captures via
%% disjoint result contexts using $\ctwand$; we elaborate on this point
%% in \autoref{sec:extension}.

\paragraph{Contributions.}
The paper makes the following contributions:\footnote{The type system,
  logic, and metatheory presented in the paper have been fully
  mechanized in Rocq, and are included in the supplemental material.}
\begin{itemize}[leftmargin=*]
\item We interpret the capability of a context to make angelic or
  demonic choices as a resource, and introduce a resource algebra to
  express compositional operations over this capability.
\item Building on this algebra, we define an associated modal
  \emph{context logic} that supports dependent and disjoint mixing of
  safety/demonic and reachability/angelic modalities.
\item We introduce \emph{context types}, a bunched modal type system
  interpreted in our context logic, to support both safety and
  reachability reasoning for deterministic functional programs.
\item We further generalize context types to a functional language
  extended with nondeterminism, where the result context is allowed to
  be disjoint from the input context.
\end{itemize}

The remainder of the paper is organized as follows.  In the next
section, we provide an example-driven overview of the type system.
Section~\ref{sec:context-typing} formalizes this informal explanation
in terms of a typed $\lambda$-calculus core language.
Section~\ref{sec:meta} defines a denotational semantics and metatheory
for this type system via a bunched implication resource capability
algebra.  Section~\ref{sec:extension} extends the language with
non-determinism, along with refinements to the logic and denotational
semantics necessary to support this new capability.
Section~\ref{sec:case} presents several case studies for which context
types are particularly well-aligned.  A brief discussion on potential
ways to realize an efficient typing algorithm as well as improved
automation for this system is given in Section~\ref{sec:discussion}.
Related work and conclusions are presented in
Sections~\ref{sec:related} and~\cite{sec:conclusion}.
